\begin{proofbox}
\textbf{证明一（优先）：沿抛物线的函数方向代入。}
对 $y^2=2spx$，由于 $s^2=1$，可以直接写成 $x=sy^2/(2p)$。
代入直线并乘以非零数 $2sp$：
\[
 A\frac{sy^2}{2p}+By+C=0
 \ \Longleftrightarrow\ Ay^2+2spBy+2spC=0.
\]
当 $A\ne0$ 时，由韦达得到纵坐标和、积。反过来，每个实根 $y$ 唯一给出 $x=sy^2/(2p)$，满足原方程组，没有增根或漏点。
如果忘记韦达，展开 $A(y-y_1)(y-y_2)$，比较一次项和常数项即可。

\textbf{配方再给一种低门槛看法。} 将二次式乘以 $A$ 并配方：
\[
 (Ay+spB)^2=p^2B^2-2spAC=H.
\]
平方非负，所以 $H<0$ 无实根；$H=0$ 只有一个重根；$H>0$ 得到两个不同实根。
当 $H>0$ 时还直接读出
\[
 y=\frac{-spB\pm\sqrt H}{A},\qquad
 |y_2-y_1|=\frac{2\sqrt H}{|A|}.
\]
$A=0$ 时不能除以 $A$，但原直线直接给出 $y=-C/B$，再求 $x$ 即可。
沿此水平线，$y^2-2spx$ 关于 $x$ 是斜率非零的一次式，经过交点时变号，直线穿过曲线，非相切。

对 $x^2=2spy$，从 $y=sx^2/(2p)$ 直接代入，得到
\[
 Bx^2+2spAx+2spC=0,\qquad (Bx+spA)^2=K.
\]
因此上下开口的韦达、判别式和轴向直线分类也全部成立。
\end{proofbox}
\clearpage
\begin{proofbox}
\textbf{证明二：用反射减少符号记忆。}
把右开口图形关于 $y$ 轴反射，设新坐标为 $(X,Y)=(-x,y)$：
\[
 y^2=2px\ \Longrightarrow\ Y^2=-2pX,\qquad
 Ax+By+C=0\ \Longrightarrow\ -AX+BY+C=0.
\]
\begin{center}\begin{tikzpicture}[font=\small]
\draw[black!30,rounded corners=2pt] (0,0) rectangle (14.4,8.2);
\node[anchor=west,font=\small\bfseries] at (.5,7.72) {左右开口互为反射：曲线、直线、交点一起变换};
% CURVE lateral 2 1
\begin{scope}[shift={(3.6,3.7)},scale=0.75]
\begin{scope}
\clip (-2.6,-2.6) rectangle (2.6,2.6);
\draw[blue,very thick,domain=-2.6:2.6,samples=181,variable=\t] plot ({\t*\t/(4)},\t);
\end{scope}
\draw[-Stealth,black!45] (-2.6,0)--(2.6,0) node[right] {$x$};
\draw[-Stealth,black!45] (0,-2.6)--(0,2.6) node[above] {$y$};
\draw[orange!85!black,very thick] (1,-2.6)--(1,2.6);
\coordinate (P) at (1.0000000000,2.0000000000);\fill[black] (P) circle (2pt) node[right=10pt] {$P$};
\coordinate (Q) at (1.0000000000,-2.0000000000);\fill[black] (Q) circle (2pt) node[right=10pt] {$Q$};
\end{scope}% CURVE lateral 2 -1
\begin{scope}[shift={(10.6,3.7)},scale=0.75]
\begin{scope}
\clip (-2.6,-2.6) rectangle (2.6,2.6);
\draw[blue,very thick,domain=-2.6:2.6,samples=181,variable=\t] plot ({\t*\t/(-4)},\t);
\end{scope}
\draw[-Stealth,black!45] (-2.6,0)--(2.6,0) node[right] {$x$};
\draw[-Stealth,black!45] (0,-2.6)--(0,2.6) node[above] {$y$};
\draw[orange!85!black,very thick] (-1,-2.6)--(-1,2.6);
\coordinate (P') at (-1.0000000000,2.0000000000);\fill[black] (P') circle (2pt) node[left=10pt] {$P'$};
\coordinate (Q') at (-1.0000000000,-2.0000000000);\fill[black] (Q') circle (2pt) node[left=10pt] {$Q'$};
\end{scope}\node at (3.6,6.7) {$y^2=4x,\ x=1$};\node at (10.6,6.7) {$y^2=-4x,\ x=-1$};
\node[anchor=west] at (.7,.55) {$(x,y)\mapsto(-x,y)$；横坐标变号，纵坐标不变，$p$ 仍然为正};
\end{tikzpicture}\end{center}

因此从向右公式得到向左公式时，需同步改变横坐标及直线中的 $A$，纵坐标保持不变。
例如原横坐标和中的 $B^2p-AC$，变换后给出向左的 $-(B^2p+AC)$；横坐标积仍为 $C^2/A^2$。
这正是第一部分的两组结果。\textbf{$p>0$ 保持不变，方向用方程中的符号表示。}
\end{proofbox}
\clearpage
\begin{proofbox}
\textbf{证明三：交换坐标，左右开口变上下开口。}
令 $(X,Y)=(y,x)$，则
\[
 y^2=2spx\ \Longrightarrow\ X^2=2spY,\qquad
 Ax+By+C=0\ \Longrightarrow\ BX+AY+C=0.
\]
\begin{center}\begin{tikzpicture}[font=\small]
\draw[black!30,rounded corners=2pt] (0,0) rectangle (14.4,8.2);
\node[anchor=west,font=\small\bfseries] at (.5,7.72) {交换坐标：右开口变上开口，竖直弦变水平弦};
% CURVE lateral 2 1
\begin{scope}[shift={(3.6,3.7)},scale=0.72]
\begin{scope}
\clip (-2.6,-2.6) rectangle (2.6,2.6);
\draw[blue,very thick,domain=-2.6:2.6,samples=181,variable=\t] plot ({\t*\t/(4)},\t);
\end{scope}
\draw[-Stealth,black!45] (-2.6,0)--(2.6,0) node[right] {$x$};
\draw[-Stealth,black!45] (0,-2.6)--(0,2.6) node[above] {$y$};
\draw[orange!85!black,very thick] (1,-2.6)--(1,2.6);
\coordinate (P) at (1.0000000000,2.0000000000);\fill[black] (P) circle (2pt) node[right=10pt] {$P$};
\coordinate (Q) at (1.0000000000,-2.0000000000);\fill[black] (Q) circle (2pt) node[right=10pt] {$Q$};
\end{scope}% CURVE vertical 2 1
\begin{scope}[shift={(10.6,3.7)},scale=0.72]
\begin{scope}
\clip (-2.6,-2.6) rectangle (2.6,2.6);
\draw[blue,very thick,domain=-2.6:2.6,samples=181,variable=\t] plot (\t,{\t*\t/(4)});
\end{scope}
\draw[-Stealth,black!45] (-2.6,0)--(2.6,0) node[right] {$x$};
\draw[-Stealth,black!45] (0,-2.6)--(0,2.6) node[above] {$y$};
\draw[orange!85!black,very thick] (-2.6,1)--(2.6,1);
\coordinate (P') at (2.0000000000,1.0000000000);\fill[black] (P') circle (2pt) node[above=12pt] {$P'$};
\coordinate (Q') at (-2.0000000000,1.0000000000);\fill[black] (Q') circle (2pt) node[above=12pt] {$Q'$};
\end{scope}\node at (3.6,6.7) {$y^2=4x$};\node at (10.6,6.7) {$x^2=4y$};
\node[anchor=west] at (.7,.55) {$(x,y)\mapsto(y,x)$，同时 $A\leftrightarrow B$；向下开口再作反射};
\end{tikzpicture}\end{center}

同时交换 $x,y$ 与 $A,B$：纵坐标二次式变成横坐标二次式，$H$ 变成 $K$，中点的两个坐标交换。
向下开口取 $s=-1$；也可以把向上图形关于 $x$ 轴反射，并同步把直线的 $B$ 变号。
因此四种开口不必分别背一套互不相关的公式。
\end{proofbox}
\clearpage
\begin{proofbox}
\textbf{横坐标固定形式为什么正确，又有什么限制？}
由直线得 $By=-(Ax+C)$，平方并代入 $y^2=2spx$：
\[
 (Ax+C)^2=2spB^2x
 \ \Longrightarrow\ A^2x^2+(2AC-2spB^2)x+C^2=0.
\]
当 $A\ne0$，韦达立即给出第一部分的左右开口横坐标关系。
\begin{center}\begin{tikzpicture}[font=\small]
\draw[black!30,rounded corners=2pt] (0,0) rectangle (14.4,8.2);
\node[anchor=west,font=\small\bfseries] at (.5,7.72) {平方后的式子相同，恢复纵坐标仍要用原直线};
% CURVE lateral 2 1
\begin{scope}[shift={(4.9,3.9)},scale=0.5]
\begin{scope}
\clip (-1,-5.5) rectangle (6.8,5.5);
\draw[blue,very thick,domain=-5.2:5.2,samples=181,variable=\t] plot ({\t*\t/(4)},\t);
\end{scope}
\draw[-Stealth,black!45] (-1,0)--(6.8,0) node[right] {$x$};
\draw[-Stealth,black!45] (0,-5.5)--(0,5.5) node[above] {$y$};
\draw[orange!85!black,very thick] (0,-1)--(6.4,5.4);
\draw[red!75!black,very thick,dashed] (0,1)--(6.4,-5.4);
\fill[black] (0.1715728753,-0.8284271247) circle (2pt);
\fill[black] (0.1715728753,0.8284271247) circle (2pt);
\fill[black] (5.8284271247,4.8284271247) circle (2pt);
\fill[black] (5.8284271247,-4.8284271247) circle (2pt);
\end{scope}\node[anchor=west,orange!85!black] at (9.5,6.5) {$y=x-1$};\node[anchor=west,red!75!black] at (9.5,5.5) {$y=1-x$};\node[anchor=west] at (9.5,3.7) {横坐标根相同};\node[anchor=west] at (9.5,2.7) {纵坐标符号相反};
\node[anchor=west] at (.7,.55) {两条直线都给出 $(x-1)^2=4x$；不能随意选 $y=\pm2\sqrt{x}$};
\end{tikzpicture}\end{center}

若 $B\ne0$，对任一实根指定 $y=-(Ax+C)/B$，由上式可反推 $y^2=2spx$，故与原交点完全对应。\textbf{恢复符号要用原直线，不能随意选正负号。}
若 $B=0$，平方后只保留 $x=-C/A$，而抛物线上的实点还需另行检验。
例如 $y^2=4x$ 与 $x=-1$：横坐标方程有重根 $-1$，原方程组却无实解。
对应的 $\Delta_x=4B^2H$ 只有在 $B\ne0$ 时，才能借其符号判断实际交点数。
\end{proofbox}
\clearpage
\begin{proofbox}
\textbf{中点短证：两点方程作差，少算一遍韦达。}
对 $y^2=2spx$ 的两个不同交点，必有 $y_1\ne y_2$。
由两点方程作差、直线方程作差，得到
\[
 (y_1-y_2)(y_1+y_2)=2sp(x_1-x_2),\qquad
 A(x_1-x_2)+B(y_1-y_2)=0.
\]
$A\ne0$ 时除以 $y_1-y_2$，立即有 $y_1+y_2=-2spB/A$。
所以 $y_S=-spB/A$，再代入 $Ax_S+By_S+C=0$，得到中点公式。
\begin{center}\begin{tikzpicture}[font=\small]
\draw[black!30,rounded corners=2pt] (0,0) rectangle (14.4,8.2);
\node[anchor=west,font=\small\bfseries] at (.5,7.72) {根差补成距离；两坐标分别取平均才是弦中点};
% CURVE lateral 2 1
\begin{scope}[shift={(3.8,5.1)},scale=0.6]
\begin{scope}
\clip (-1,-6.2) rectangle (8.3,2.2);
\draw[blue,very thick,domain=-6.0:2.2,samples=181,variable=\t] plot ({\t*\t/(4)},\t);
\end{scope}
\draw[-Stealth,black!45] (-1,0)--(8.3,0) node[right] {$x$};
\draw[-Stealth,black!45] (0,-6.2)--(0,2.2) node[above] {$y$};
\coordinate (P) at (0.5358983849,1.4641016151);\fill[black] (P) circle (2pt) node[above left=9pt,xshift=-16pt] {$P$};
\coordinate (Q) at (7.4641016151,-5.4641016151);\fill[black] (Q) circle (2pt) node[above right=11pt] {$Q$};
\coordinate (T) at (7.4641016151,1.4641016151);\fill[black] (T) circle (2pt) node[above right=9pt] {$T$};
\draw[orange!85!black,very thick] (0.5358983848622456,1.4641016151377544)--(7.464101615137754,-5.464101615137754);
\draw[teal!70!black,dashed,thick] (P)--(T)--(Q);
\draw (7.1641016151,1.4641016151)--(7.1641016151,1.1641016151)--(7.4641016151,1.1641016151);
\coordinate (S) at (4.0000000000,-2.0000000000);\fill[violet] (S) circle (2pt) node[above right=10pt] {$S$};
\node[above=12pt] at (4,1.4641016151) {$|x_2-x_1|$};\node[anchor=west] at (8.1,-2) {$|y_2-y_1|$};\end{scope}
\node[anchor=west] at (.7,.55) {$PQ^2=(x_2-x_1)^2+(y_2-y_1)^2$；$S=(4,-2)$ 在弦上，不在抛物线上};
\end{tikzpicture}\end{center}

\textbf{弦长只需补一个直角三角形。}
由差分关系 $x_2-x_1=-(B/A)(y_2-y_1)$，勾股定理给出
\[
 PQ=\frac{\sqrt{A^2+B^2}}{|A|}\,|y_2-y_1|
 =\frac{2\sqrt{(A^2+B^2)H}}{A^2}.
\]
这里允许 $B=0$，所以竖直弦也已覆盖。上下开口交换坐标即得对应中点和弦长公式。
\end{proofbox}
\clearpage
